% English
\usepackage[english]{babel} % English language setting
\usepackage[utf8]{inputenc} % Unicode text
\usepackage[T1]{fontenc} % German 'Umlaute'
\usepackage{textcomp} % Euro
\usepackage[hyphens]{url}
\usepackage{amsmath} % Advanced mathematical notation
% \usepackage{amssymb} % Symbols
\usepackage{emptypage} % Empty pages are now actually empty
\usepackage{caption} % Better image and table captions
\usepackage{varwidth} % Used for the custom captions
\DeclareCaptionFormat{centered}{%
  % #1: label (e.g. "Table 1")
  % #2: separator (e.g. ": ")
  % #3: caption text
  \begin{varwidth}{\linewidth}
    #1#2#3%
  \end{varwidth}
}
\captionsetup{format=centered} % Global activation of the custom captions

\usepackage{color}
\usepackage[dvipsnames]{xcolor}
\usepackage[tikz]{bclogo}
\usepackage[framemethod=tikz]{mdframed}
\usepackage{lipsum}
\usepackage[many]{tcolorbox}
\newtcolorbox{highlightbox}[1][]{
  breakable,
  freelance,
  title=#1,
  colback=white,
  colbacktitle=white,
  coltitle=black,
  fonttitle=\bfseries,
  bottomrule=0pt,
  boxrule=0pt,
  colframe=white,
  overlay unbroken and first={
  \draw[black,line width=3pt]
    ([xshift=5pt]frame.north west) --
    (frame.north west) --
    (frame.south west);
  \draw[black,line width=3pt]
    ([xshift=-5pt]frame.north east) --
    (frame.north east) --
    (frame.south east);
  },
  overlay unbroken app={
  \draw[black,line width=3pt,line cap=rect]
    (frame.south west) --
    ([xshift=5pt]frame.south west);
  \draw[black,line width=3pt,line cap=rect]
    (frame.south east) --
    ([xshift=-5pt]frame.south east);
  },
  overlay middle and last={
  \draw!black,line width=3pt]
    (frame.north west) --
    (frame.south west);
  \draw[black,line width=3pt]
    (frame.north east) --
    (frame.south east);
  },
  overlay last app={
  \draw[black,line width=3pt,line cap=rect]
    (frame.south west) --
    ([xshift=5pt]frame.south west);
  \draw[black,line width=3pt,line cap=rect]
    (frame.south east) --
    ([xshift=-5pt]frame.south east);
  },
}

% Fonts, with all the options
\usepackage[bitstream-charter]{mathdesign}
\usepackage[scaled=.96]{helvet}
\usepackage[scaled=.75]{beramono}
\usepackage{microtype}

% Images and listings
\usepackage{graphicx} % images
\usepackage{subfig} % sub-figures
\usepackage{wrapfig} % wrapping figures
\usepackage{listings} % better source code listings
\usepackage[algoruled,lined,linesnumbered,commentsnumbered]{algorithm2e} % algorithms
\usepackage{tabularx} % better tables
\usepackage{diagbox} % diagonally split table cells
\usepackage{multirow}
\usepackage{rotating}
\usepackage{hhline}
\usepackage{enumitem}
\usepackage{array}
\usepackage{booktabs}
\usepackage{url}
\usepackage{forest}
\usepackage{verbatim}
\usepackage{pgfplots}
\pgfplotsset{width=7cm,compat=1.8}
\usepackage{pgfplotstable}
\graphicspath{ {figures/} }

% Credit: https://github.com/cansik/kotlin-latex-listing
\lstdefinelanguage{Kotlin}{
  comment=[l]{//},
  commentstyle={\color{gray}\ttfamily},
  emph={delegate, filter, first, firstOrNull, forEach, lazy, map, mapNotNull, println, return@},
  emphstyle={\color{OrangeRed}},
  identifierstyle=\color{black},
  keywords={abstract, actual, as, as?, break, by, class, companion, continue, data, do, dynamic, else, enum, expect, false, final, for, fun, get, if, import, in, interface, internal, is, null, object, override, package, private, public, return, set, super, suspend, this, throw, true, try, typealias, val, var, vararg, when, where, while},
  keywordstyle={\color{NavyBlue}\bfseries},
  morecomment=[s]{/*}{*/},
  morestring=[b]",
  morestring=[s]{"""*}{*"""},
  ndkeywords={@Deprecated, @JvmField, @JvmName, @JvmOverloads, @JvmStatic, @JvmSynthetic, Array, Byte, Double, Float, Int, Integer, Iterable, Long, Runnable, Short, String},
  ndkeywordstyle={\color{BurntOrange}\bfseries},
  sensitive=true,
  stringstyle={\color{ForestGreen}\ttfamily},
}

\definecolor{lstgrey}{rgb}{0.95,0.95,0.95}
\lstset{
    language=Kotlin,
    frame=single,
    backgroundcolor=\color{lstgrey},
    basicstyle=\small\ttfamily,
    captionpos=b,
    tabsize=2,
    commentstyle=\color{gray},              % Comments font
    tabsize=2,                              % Default tab size
    captionpos=b,                           % Caption-position = bottom
    breaklines=true,                        % Automatic line breaking?
    breakatwhitespace=false,                % Automatic breaks only at whitespace?
    showspaces=false,                       % Dont make spaces visible
    showtabs=false,                         % Dont make tabs visible
    columns=flexible,                       % Column format
}

% Page layout
% Comment following two lines in for proper layout
% \usepackage[paper=a4paper,width=14cm,left=35mm,height=22cm]{geometry}
% \linespread{1.5}
\usepackage[paper=a4paper]{geometry}
\usepackage{setspace}
\usepackage[htt]{hyphenat}
\usepackage{sectsty}
\subsubsectionfont{\large}

% Page markers
\newcommand{\phv}{\fontfamily{phv}\fontseries{m}\fontsize{10}{12}\selectfont}
\usepackage{fancyhdr} % nicer header and footer
\pagestyle{fancy}
\renewcommand{\chaptermark}[1]{\markboth{#1}{}}
\fancyhead[L]{\phv \nouppercase{\leftmark}}
\fancyhead[R]{\phv \thepage}
% rather not use anything for the footer
\fancyfoot[C]{\ } % no page count in the bottom
% \fancyfoot[R]{\textsf{\small Software Engineering MS}}

% Share the sources
\usepackage{bibtopic}

% Special packages
\usepackage{epigraph}
\setlength{\epigraphrule}{0pt} % no divider
\usepackage{csquotes}

% Some extra styles
\usepackage{soul}
\newcommand*\strikethrough{\st}

% Hyperlink everything
\usepackage{hyperref}
\hypersetup{
    colorlinks=true,
    linkcolor=black,
    citecolor=black,
    filecolor=black,
    urlcolor=black,
}
\urlstyle{same}

% Hyphenations
\hyphenation{non-neg-li-gi-ble}

\newcommand*\rot{\rotatebox{90}}
\newcolumntype{Y}{>{\centering\arraybackslash}X}

% Wikipedia-style "citation needed" macro
\newcommand{\cn}[1][]{\textsuperscript{\color{red} ~[citation needed]~}}

% UvA logo and text
\newcommand{\uva}{
	\begin{tabular}{r l}
		\multirow{2}{*}{\includegraphics[height=55pt]{uva.pdf}}
		&\textsc{\Large Universiteit van Amsterdam}\\
		&\textsc{\small Faculteit der Natuurwetenschappen, Wiskunde en Informatica}\\
		&\textsc{\small Master Software Engineering}
	\end{tabular}
}

% Code markup
\newcommand{\code}[1]{\texttt{#1}}
